\chapter{The device library}
\label{ch:library}

The Library page keeps your device maps in one folder, and brings maps from
the shared online library.

\screenshot[7cm]{library}{The Library page.}
  {The Library page with several maps: vendor icons, vendor, model,
   description, number of registers and status icons; the search field and the
   two import buttons.}

\section{Your library}

The list shows each map's vendor (with its icon), model, description, number
of registers and status (Section~\ref{sec:map-status}). \field{Search vendor or
model} filters it.

Right-click a map:

\begin{description}[style=nextline, leftmargin=0pt, font=\intermedium]
  \item[Use] Opens it on the Main page.
  \item[Simulate this device] Opens the simulator with a device shaped after
    this map (Chapter~\ref{ch:simulator}).
  \item[Copy name, Copy file path] For notes and other programs.
  \item[Open the library folder] Shows the folder in your file manager.
  \item[Delete…] Removes the map from the library, after asking.
\end{description}

Maps with errors are listed with the problem, so they can be fixed in a text
editor.

\section{Adding maps}

\begin{description}[style=nextline, leftmargin=0pt, font=\intermedium]
  \item[\ui{From a file…}] Copies \file{.mbmap.json} or \file{.mbmap} files into
    the library. Version 1 files are converted.
  \item[\ui{From the online library…}] Opens the shared library (below).
\end{description}

The library is a folder (\nav{Settings\then Device library\then Library
folder}): files copied into it by hand appear too.

\section{The online library}

\screenshot[7cm]{online-library}{Choosing a map in the online library.}
  {The online library: vendors on the left with their icons, the models of one
   vendor on the right, one map selected.}

The online library lists the maps shared on GitHub
(\web{github.com/nobrega8/mbmap-library}): over 400 maps of meters,
inverters, drives and controllers.

\begin{steps}
  \item Choose a vendor; \ui{Back to all vendors} (or \key{Esc}) goes back.
  \item \ui{Add} next to a model downloads its map into your library.
  \item \ui{Done} closes the online library.
\end{steps}

A map already in your library can be downloaded again, to replace your copy
with the newest one.

\begin{tip}[Share your maps]
A map that works with a real device helps the next person with the same
device. The library's repository explains how to contribute one, and the
format specification is linked from \nav{Settings\then Device library}.
\end{tip}
